\section{Ranks and decomposition certificates}\label{sec:preliminaries}

All coordinate and record sets are finite. A \emph{polymatroid} is a
function $h:2^A\to\R$ with $h(\varnothing)=0$, monotonicity under
inclusion, and
\[
h(S\cup T)+h(S\cap T)\le h(S)+h(T).
\]
Write
\begin{align*}
h(S\mid K)&=h(S\cup K)-h(K),\\
I_h(S;T\mid Z)&=h(S\mid Z)-h(S\mid T\cup Z),\\
h_Z(S)&=h(S\mid Z).
\end{align*}
A singleton may be written without braces. Contraction $h_Z$ and
pullback along any map $f:V\to A$ are again polymatroids. The map need
not be injective. These are the existing contraction theorem and the
memo's \lean{image_isPolymatroid} and
\lean{conditionalImage_isPolymatroid}.

For jointly distributed finite random variables $W_a$, Shannon entropy
$h(S)=H(W_S)$ supplies an instance. We use natural logarithms, as does
the Lean entropy library: a fair bit has entropy $\log 2$. All terms
must use the same units. Symbols such as $H(L\mid K,Z)$ abbreviate the
entropy of the corresponding tuples. No independence is implicit.

\paragraph{The imported structural certificate.}
Let $V$ be a finite vertex set and $N:V\to2^V$ a prescribed-parent
relation. We recall the certificate of \cite{SeventhPaper2026} to make
clear which hypothesis is retained. A piece $H\subseteq V\setminus U$
has a certificate $\PieceCert(N,U,H)$ consisting of:
\begin{enumerate}[label=(\roman*)]
\item finitely many pairwise disjoint sets $P_1,\ldots,P_m$ whose union
is $H$;
\item a list $\ell$ whose set of entries is $H$, with repetitions
allowed, and maps $a,b:V\to\{1,\ldots,m\}$ such that
$v\in P_{a(v)}$ and $a(v)\ne b(v)$ for every entry;
\item a fixed optional outside parent $\operatorname{par}(v)$ for each
vertex, satisfying the following scan conditions.
\end{enumerate}
The scan starts with $S=\varnothing$ and $M_w=\varnothing$ for every
$w\in V$. At an entry $v$, put $K=S\cup P_{b(v)}$. If there is no
outside parent, require $N(v)\subseteq U\cup K$. If the outside parent
is $w$, require
\[
w\in V\setminus(U\cup H),\qquad
N(v)\subseteq U\cup K\cup\{w\},\qquad M_w\subseteq K,
\]
and update $M_w$ to $K\cup\{v\}$. In both cases insert $v$ into $S$.
The supplying part $P_{b(v)}$ is part of this structural certificate;
it is distinct from the actual helper records introduced later.

An ordered list of pieces $\mathcal H=(H_1,\ldots,H_t)$ satisfies
$\Decomp(N,U,\mathcal H)$ when the pieces partition $V\setminus U$
and each $H_j$ has a piece certificate in context
\[
U_j=U\cup\bigcup_{i<j}H_i.
\]
This is precisely the imported predicate
\leanfile{HalfRate.Decomp}, with piece type
\leanfile{HalfRate.PieceCert}. These definitions contain no rank
function or recovery equality. They are not inferred merely from the
chronological order of a tape computation.

\paragraph{Class S.}
An arc $u\to v$ means $u\in N(v)$. Class S consists of loopless
relations with no transitive triangle on three distinct vertices,
in-degree at most three, and the property that every vertex of
in-degree three reciprocates each outgoing arc. The relative condition
on $R\subseteq V$ imposes these requirements on the induced relation
inside $R$, counting only parents in $R$. The tenth paper proves that
this relative condition supplies
$\Decomp(N,V\setminus R,\mathcal H)$ for some $\mathcal H$
\cite[Theorem S*]{TenthPaper2026}; its proof is Lean-verified there.
